\begin{table}[t]
\centering
\caption{Event Study: Year-Specific Effects on Gender Wage Ratio}
\label{tab:eventstudy}
\begin{tabular}{lcc}
\hline\hline
Year $\times$ Treatment Intensity & Coefficient & Std. Error \\
\hline
2014 \textit{(pre)} & $5.86$ & $(9.75)$ \\
2015 \textit{(pre)} & $-4.57$ & $(6.99)$ \\
2017 \textit{(post)} & $0.99$ & $(7.37)$ \\
2018 \textit{(post)} & $4.04$ & $(7.12)$ \\
2019 \textit{(post)} & $-4.58$ & $(6.42)$ \\
2020 \textit{(post)} & $-3.58$ & $(6.42)$ \\
2021 \textit{(post)} & $7.91$ & $(5.82)$ \\
2022 \textit{(post)} & $16.29^{**}$ & $(7.34)$ \\
2023 \textit{(post)} & $12.18$ & $(10.18)$ \\
2024 \textit{(post)} & $11.56$ & $(9.79)$ \\
\hline
Reference year & \multicolumn{2}{c}{2016} \\
Observations & \multicolumn{2}{c}{209} \\
Industry FE & \multicolumn{2}{c}{Yes} \\
Year FE & \multicolumn{2}{c}{Yes} \\
\hline\hline
\multicolumn{3}{p{0.85\textwidth}}{\footnotesize \textit{Notes:} Each row reports the coefficient on $\mathbf{1}[\text{Year}=t] \times \text{TreatIntensity}_{i}$, with 2016 as the omitted reference year. The outcome is the gender wage ratio (female salary as \% of male). Standard errors clustered by industry. $^{*}p<0.10$, $^{**}p<0.05$, $^{***}p<0.01$.} \\
\end{tabular}
\end{table}
